% =============================================================================================
%  Recherche de la section du C-foil par co-krigeage séquentiel ciblé par la VPP :
%  description exacte de l'algorithme et de l'architecture (état du code du 26/10/2026,
%  branche feat/vpp-acceleration, 4e vrai run ; première version du 22/10, 3e run).
%  Compilation : tectonic -X compile algorithme_vpp_cible.tex --outdir build
% =============================================================================================
\documentclass[10pt,a4paper]{article}
\usepackage{fontspec}
\usepackage[french]{babel}
\usepackage{csquotes}
\usepackage[a4paper,margin=1.8cm,top=1.6cm,bottom=1.8cm]{geometry}
\usepackage{microtype}
\usepackage{amsmath,amssymb,amsthm,mathtools,bm}
\usepackage{xcolor}
\usepackage{booktabs,tabularx,array,longtable}
\usepackage{enumitem}
\usepackage[most]{tcolorbox}
\usepackage{listings}
\usepackage{algorithm}
\usepackage{algpseudocode}
\usepackage{titlesec}
\usepackage{tikz}
\usetikzlibrary{arrows.meta,positioning,fit,backgrounds}
\usepackage[hidelinks,unicode]{hyperref}

\hypersetup{colorlinks=true,linkcolor=blue!50!black,citecolor=green!40!black,
            urlcolor=blue!60!black,pdftitle={Recherche de la section du C-foil ciblee par la VPP}}
\setlength{\parindent}{0pt}
\setlength{\parskip}{0.4em}
\setlist{itemsep=0.1em,topsep=0.2em,leftmargin=1.3em}
\titleformat{\section}{\normalfont\large\bfseries\color{blue!40!black}}{\thesection}{0.6em}{}
\titlespacing*{\section}{0pt}{0.9em}{0.3em}
\titleformat{\subsection}{\normalfont\normalsize\bfseries}{\thesubsection}{0.5em}{}
\titlespacing*{\subsection}{0pt}{0.6em}{0.2em}
\renewcommand{\arraystretch}{1.12}

\tcbset{enhanced,breakable,boxrule=0.6pt,arc=1.2mm,left=1.5mm,right=1.5mm,top=0.6mm,bottom=0.6mm,
        fonttitle=\bfseries\small,before skip=0.5em,after skip=0.5em}
\newtcolorbox{danslecode}[1][]{colback=green!3,colframe=green!35!black,
  title={Dans le code\ifx&#1&\else~— #1\fi}}
\newtcolorbox{subtilite}[1][]{colback=orange!4,colframe=orange!70!black,
  title={Subtilité\ifx&#1&\else~— #1\fi}}
\newtcolorbox{aretenir}[1][]{colback=violet!4,colframe=violet!60!black,
  title={À retenir\ifx&#1&\else~— #1\fi}}
\newtcolorbox{nonfait}[1][]{colback=red!3,colframe=red!50!black,
  title={Ce qui n'est pas fait\ifx&#1&\else~— #1\fi}}

\floatname{algorithm}{Algorithme}
\algrenewcommand\algorithmicrequire{\textbf{Entrées :}}
\algrenewcommand\algorithmicensure{\textbf{Sorties :}}
\algrenewcommand\algorithmicwhile{\textbf{tant que}}
\algrenewcommand\algorithmicdo{\textbf{faire}}
\algrenewcommand\algorithmicfor{\textbf{pour}}
\algrenewcommand\algorithmicif{\textbf{si}}
\algrenewcommand\algorithmicthen{\textbf{alors}}
\algrenewcommand\algorithmicelse{\textbf{sinon}}
\algrenewcommand\algorithmicend{\textbf{fin}}
\algrenewcommand\algorithmicreturn{\textbf{retourner}}

\newcommand{\vx}{\bm{x}}
\newcommand{\vz}{\bm{z}}
\newcommand{\vw}{\bm{w}}
\newcommand{\va}{\bm{a}}
\newcommand{\vu}{\bm{u}}
\newcommand{\vd}{\bm{d}}
\newcommand{\vc}{\bm{c}}
\newcommand{\D}{\mathcal{D}}
\newcommand{\tp}{^{\mathsf{T}}}
\DeclareMathOperator{\Var}{Var}
\DeclareMathOperator{\Cov}{Cov}
\DeclareMathOperator*{\argmax}{arg\,max}
\ExplSyntaxOn
\cs_new_protected:Npn \mf_breakable:n #1
  {
    \str_set:Nn \l_tmpa_str {#1}
    \str_map_inline:Nn \l_tmpa_str
      {
        ##1
        \str_if_in:nnT { _./,(=[ } {##1} { \allowbreak }
      }
  }
\NewDocumentCommand \py { m } { \texttt { \mf_breakable:n {#1} } }
\ExplSyntaxOff

\begin{document}

{\centering
{\Large\bfseries\color{blue!40!black} Recherche de la section du C-foil par co-krigeage séquentiel
ciblé par la VPP\par}
\vspace{0.2em}
{\large Description exacte de l'algorithme et de l'architecture\par}
\vspace{0.3em}
{\small \py{Multifidelity-GP-Clean/vpp_acceleration} · branche \py{feat/vpp-acceleration} · état du
code du 4e vrai run (\py{runs/2026-10-23_1702_real_search_profile_local}) · 26/10/2026\par}
\vspace{0.2em}}

\tableofcontents

% ---------------------------------------------------------------------------------------------
\section{Ce que fait l'algorithme, en une page}
% ---------------------------------------------------------------------------------------------
\textbf{Le but.} On cherche la section (le profil) du C-foil de l'Intens SY qui maximise la vitesse du
bateau donnée par la VPP de MDS, moyennée sur des conditions de vent pondérées par leur probabilité : cinq
conditions de même poids aux trois premiers runs, les 18 nœuds d'une distribution d'usage au
4\textsuperscript{e}. La forme en plan du foil est fixée, et seule la section change.

\textbf{La difficulté.} La VPP appelle le modèle du foil des centaines de fois par point de polaire. Les codes
aérodynamiques (NPLLT, AVL avec XFOIL) sont trop lents pour cela. On remplace donc le foil, dans la VPP, par
un \textbf{surrogate} de son torseur : un co-krigeage multi-fidélité, dont les entrées sont la section et
l'attitude du foil. Ce surrogate s'améliore pendant la recherche. On ajoute des calculs de code
\textbf{là où la VPP travaille} (les états d'équilibre des sections candidates) et \textbf{sur ce qui compte
pour la vitesse} (chaque coefficient est pondéré par $(\partial V/\partial C_k)^2$).

\textbf{La boucle.} Elle alterne deux étapes jusqu'à une limite de temps.
\begin{enumerate}
  \item \textbf{La mise à jour de l'optimiseur.}
        \begin{itemize}
          \item On crible 256 sections nouvelles au premier ordre, avec le surrogate et les sensibilités de
                vitesse de la meilleure section du moment (le \enquote{titulaire}). Depuis le 23/10, la
                moitié est tirée dans une région de confiance autour du titulaire (recherche locale).
          \item On résout la VPP complète sur les 3 meilleures et sur le titulaire, puis le titulaire devient
                la meilleure des quatre.
          \item Les états d'équilibre de ces 4 sections, à chaque condition, donnent les points
                d'intégration du critère de raffinement.
        \end{itemize}
  \item \textbf{Un lot de raffinement.} On choisit 16 couples (point, niveau de code), l'un après l'autre,
        par la meilleure réduction de variance de la vitesse par unité de coût. Puis les codes du lot
        tournent en parallèle, et leurs résultats entrent dans le co-krigeage.
\end{enumerate}

\textbf{Les niveaux de code} (3\textsuperscript{e} et 4\textsuperscript{e} runs).
\begin{itemize}
  \item \textbf{Niveau 1 (BF) :} ligne portante non plane NPLLT avec les polaires NeuralFoil xxxlarge.
  \item \textbf{Niveau 2 (HF) :} la même ligne portante, avec les polaires XFOIL.
  \item \textbf{Niveau 3 :} AVL × XFOIL, la chaîne de la RS de production. Il ne sert qu'à la validation finale.
\end{itemize}

\textbf{Le foil est flexible.} Le surrogate et les codes sont rigides. La VPP ajoute une correction de
déformée, tirée de la RS FSI de production et mise à l'échelle de la raideur de la section (§\,\ref{sec:fsi}).

\begin{figure}[h]
\centering
\begin{tikzpicture}[font=\footnotesize, >=Stealth,
  box/.style={draw, rounded corners=2pt, align=center, fill=white, inner sep=3pt, minimum height=0.9cm},
  pool/.style={draw, rounded corners=2pt, align=center, fill=orange!8, inner sep=3pt, minimum height=0.9cm},
  disk/.style={draw, align=center, fill=gray!10, inner sep=3pt}]
  \node[box, text width=4.3cm] (loop) {\textbf{Boucle} \py{VPPSectionOptimizer}\\ mises à jour, lots, arrêt au temps limite};
  \node[box, text width=4.3cm, below=0.35cm of loop] (gp) {\textbf{Surrogate} \py{TorsorCoKriging}\\ 7 \py{CoKriging} (Fx … Mz, Cpmin), 7 fils};
  \node[box, text width=4.3cm, below=0.35cm of gp] (crit) {\textbf{Critère} \py{choose} / \py{batched_imse_gains}\\ lot de 16 par \enquote{kriging believer}};
  \begin{scope}[on background layer]
    \node[draw, dashed, rounded corners, fill=blue!3, fit=(loop)(gp)(crit), inner sep=6pt,
          label={[font=\footnotesize\bfseries]above:processus principal}] (main) {};
  \end{scope}
  \node[pool, text width=3.6cm, left=2.3cm of gp] (codes) {\textbf{\py{CodeRunner}} : 28 processus\\ \py{FoilCodes} : NPLLT × NF,\\ NPLLT × XFOIL,\\ AVL × XFOIL (bdFoil)};
  \node[disk, text width=3.2cm, below=0.5cm of codes] (cache) {cache des polaires XFOIL\\ (disque, partagé)};
  \node[pool, text width=3.6cm, right=2.3cm of gp] (vpp) {\textbf{\py{VPPPool}} : 30 processus\\ \py{IntensSYVPP} (MDS)\\ + \py{FSICorrection}};
  \node[disk, text width=3.2cm, below=0.5cm of vpp] (pkl) {\py{surrogate_v<n>.pkl}\\ (moyennes seules)};
  \draw[->] (crit.west) -- ++(-0.45,0) |- ([yshift=-2mm]codes.east) node[pos=0.72, below, font=\scriptsize] {points, niveaux};
  \draw[->] ([yshift=2mm]codes.east) -- ([yshift=2mm]gp.west) node[midway, above] {coefficients};
  \draw[<->] (codes) -- (cache);
  \draw[->] (gp.east) -- (vpp.west) node[midway, above] {publication};
  \draw[->] (vpp) -- (pkl);
  \draw[->] (vpp.north) |- (loop.east) node[pos=0.75, above] {états, $\partial V/\partial C$};
\end{tikzpicture}
\caption{Les trois groupes de processus d'un run (§\,\ref{sec:architecture}). Les deux pools ne travaillent
jamais en même temps : la boucle est synchrone.}
\label{fig:architecture}
\end{figure}

\begin{aretenir}
\begin{itemize}
  \item \textbf{Ce qui est optimisé :} 10 scores PCA d'une section de Kulfan (16 poids), dans une famille de
        sections admissibles.
  \item \textbf{Ce que voit la VPP :} la moyenne du co-krigeage, plus la correction FSI.
  \item \textbf{Ce qui est raffiné :} le co-krigeage, autour des états VPP des sections en compétition.
  \item \textbf{Ce qui décide :} la vitesse moyenne pondérée sur les conditions, calculée par la VPP complète
        sur 4 sections à chaque mise à jour.
  \item \textbf{La distribution d'usage} (TWA, TWS) se donne à chaque run (§\,\ref{sec:profil}). Elle sert à
        analyser les sections retenues (\py{--profile}) et, depuis le 23/10, peut aussi remplacer les
        conditions de la recherche (\py{--search-profile}).
\end{itemize}
\end{aretenir}

\subsection*{Notations}
\begin{tabularx}{\linewidth}{@{}lX@{}}
\toprule
$\vz\in\mathbb{R}^{k}$, $k=10$ & scores PCA de la section (blanchis, de variance 1 sur la famille) \\
$\vw = \mathcal{D}(\vz)\in\mathbb{R}^{16}$ & poids de Kulfan décodés (8 extrados, 8 intrados) \\
$\va=(a_c,a_r,a_y)$ & attitude du foil : \py{Cant_Eff}, \py{Rake_Eff}, \py{Yaw_Eff} [°] \\
$\vx=(\vz,\va)\in\mathbb{R}^{13}$ & entrée du surrogate \\
$C_k(\vx)$, $k=1..7$ & sorties : $F_x,F_y,F_z,M_x,M_y,M_z$ (coefficients, moments en m) et $C_{p,\min}$ \\
$f_l$, $l=1,2,3$ & codes : NPLLT × NeuralFoil, NPLLT × XFOIL, AVL × XFOIL \\
$\mu_k(\vx)$, $s_k^2(\vx)$ & moyenne et variance a posteriori du co-krigeage de la sortie $k$ (niveau 2) \\
$c=(\mathrm{TWA}_c,\mathrm{TWS}_c)$, $w_c$ & condition de navigation, $c=1..n_c$, et sa probabilité ($\sum_c w_c=1$) \\
$\vu\in\mathbb{R}^7$ & inconnues de la VPP : vitesse, dérive, gîte, assiette, déplacement de carène, safran,
réduction de voilure (\py{Flat}) \\
$V^\star(\vz;c)$ & vitesse d'équilibre maximale de la section $\vz$ à la condition $c$ \\
$J(\vz)$ & objectif : $\sum_c w_c V^\star(\vz;c)\,/\sum_c w_c$, sur les conditions résolues \\
$d_{kc}$ & sensibilité $\partial V^\star/\partial C_k$ à l'état du titulaire, condition $c$ [nd par unité] \\
\bottomrule
\end{tabularx}

% ---------------------------------------------------------------------------------------------
\section{Le problème d'optimisation}
% ---------------------------------------------------------------------------------------------
\[
  \max_{\vz}\; J(\vz) = \frac{\sum_{c\in\mathcal{S}} w_c\,V^\star(\vz;c)}{\sum_{c\in\mathcal{S}} w_c}
  \qquad\text{sous la contrainte}\qquad \mathcal{D}(\vz)\in\mathcal{F},
\]
où $\mathcal{F}$ est la famille des sections admissibles (§\,\ref{sec:section}), et $\mathcal{S}$ l'ensemble
des conditions où la VPP trouve un équilibre pour \emph{toutes} les sections comparées : on compare toujours
les sections sur les mêmes conditions.

\begin{itemize}
  \item \textbf{Les conditions, deux cas.}
        \begin{itemize}
          \item Par défaut (runs 1 à 3), ce sont les couples (TWA, TWS) : (45°, 12 nd), (45°, 16 nd),
                (60°, 16 nd), (75°, 16 nd) et (90°, 14 nd), de même poids $w_c=\frac15$.
          \item Avec \py{--search-profile} (4\textsuperscript{e} run), ce sont les nœuds de la distribution
                d'usage, avec leurs poids (§\,\ref{sec:profil}). Avec \py{profiles/guess_bimodal.json}, cela
                fait 18 nœuds, tous de poids $\frac1{18}$ : la probabilité est portée par leur position.
        \end{itemize}
  \item \textbf{La vitesse $V^\star(\vz;c)$} est celle de la VPP de MDS (§\,\ref{sec:vpp}), dont le foil a pour
        coefficients $\mu_k(\vz,\va)$ plus la correction FSI.
  \item \textbf{Le référentiel du gain :} la recherche part du profil actuel du C-foil, admissible une fois
        projeté sur les 10 modes. Le gain est donc mesuré par rapport au foil actuel.
\end{itemize}

\begin{danslecode}
\py{scripts/run_real_vpp.py} : \py{CONDITIONS}, ou les nœuds de \py{OperatingProfile.nodes()}, et
\py{condition_weights}. Puis \py{VPPSectionOptimizer.update} : les poids sont mis à zéro aux conditions non
résolues pour l'une des 4 sections (\py{weights = condition_weights * solved.all(axis=0)}), puis la vitesse
de chaque section est la moyenne pondérée sur les conditions restantes.
\end{danslecode}

% ---------------------------------------------------------------------------------------------
\section{La VPP : MDS, Intens SY, variante \texorpdfstring{\py{CFoilFSI}}{CFoilFSI}}
\label{sec:vpp}
% ---------------------------------------------------------------------------------------------
\subsection{Le système}
La VPP est celle de MDS, branche \py{new_master}, avec le bateau \py{build_IntensSY(foil='CFoilFSI')}
(C-foil de rayon 9,8 m). Elle reproduit vpp36 à $2\cdot10^{-6}$ nd près. C'est un \py{IntensSYSystem} de
9 sous-systèmes, évalués dans cet ordre (lu sur le système construit) :

\begin{tabularx}{\linewidth}{@{}l>{\raggedright\arraybackslash}X>{\raggedright\arraybackslash}X@{}}
\toprule
Sous-système & Entrées (variables globales) & Modèle \\
\midrule
\py{Weight} & masse, $X/Y/Z_{CG}$, gîte, assiette, dérive & poids au CG (pas de RS) \\
\py{Hull Dyn} & gîte, assiette, dérive, \py{CC_Dspl}, vitesse & RS de Chebyshev (efforts hydrodynamiques de coque) \\
\py{Hull Archi} & gîte, assiette, dérive, \py{CC_Dspl} & RS (hydrostatique, et l'enfoncement \py{CC_Sink_Out}) \\
\py{Foil Dyn} & gîte, dérive, assiette, \py{Foil_Cant/Rake/Yaw}, vitesse & \textbf{RS du C-foil, remplacée par notre surrogate} \\
\py{Foil Archi} & idem sans la vitesse & RS (poussée d'Archimède du foil) \\
\py{Rudder}, \py{Rudder STBD} & gîte, assiette, dérive, angle de safran, vitesse, \py{Hvol} & RS des deux safrans \\
\py{Sail} & dérive, \py{Flat}, TWA, TWS, vitesse, gîte, assiette, quête du mât & RS des voiles (entrées AWA, \py{Flat}) \\
\py{Windage} & gîte, assiette, dérive, TWA, TWS, vitesse & RS du fardage \\
\bottomrule
\end{tabularx}

\textbf{Les équations d'équilibre.} \py{TSystem.t_eval} évalue chaque sous-système. Chacun rend son torseur
(forces et moments) transporté au repère de référence \py{CG ref}, situé en $(X_{CG},0,0)$. La somme donne
les 6 résidus \py{Eq.Fx} … \py{Eq.Mz}. Avant cela, \py{IntensSYSystem.update_system_frame} fait deux
opérations :
\begin{itemize}
  \item il ferme la hauteur d'assise \py{Hvol}, à partir de l'enfoncement statique et dynamique de la coque ;
  \item il tourne les bras de levier (les positions des repères des composants) selon l'attitude du bateau.
        Les forces, elles, sont déjà dans le repère monde.
\end{itemize}

\subsection{Le problème résolu à chaque condition}
Les 7 inconnues $\vu$ et leurs bornes, pour la RS de coque \py{lin85t} :

\begin{tabular}{@{}llll@{}}
\toprule
Inconnue & Sens & Bornes & Départ \\
\midrule
\py{Speed} & vitesse du bateau [nd] & $[6, 24]$ & 12 \\
\py{Leeway} & dérive [°] & $[-6, 6]$ & 3 \\
\py{Heel} & gîte [°] & $[0, 8]$ & 3 \\
\py{Trim} & assiette [°] & $[-1{,}5;\ 1{,}5]$ & 0 \\
\py{CC_Dspl} & déplacement porté par la carène [kg] & $[8500, 13500]$ & 11\,540 \\
\py{Rudder_Angle} & angle de safran [°] & $[-12, 12]$ & 0 \\
\py{Flat} & réduction de puissance des voiles [–] & $[0{,}3;\ 1]$ & 0,6 \\
\bottomrule
\end{tabular}

\[
  V^\star(\vz;c) = \max_{\vu}\; u_1 \quad\text{s.c.}\quad
  \frac{R_i(\vu;\vz,c)}{e_i} = 0,\; i=1..6, \qquad \vu_{\min}\le\vu\le\vu_{\max},
\]
avec $R_i$ les résidus \py{Eq.*} et $e=(10^3,10^3,10^3,10^4,10^4,10^4)$ (N et N·m mis à la même échelle).
\begin{itemize}
  \item \textbf{Le solveur :} SLSQP de SciPy (\py{ftol} $10^{-9}$, 400 itérations au plus).
  \item \textbf{Le succès :} il faut que SLSQP converge \emph{et} que $\|R/e\|<10^{-4}$.
  \item \textbf{La relance :} en cas d'échec, une seule. Elle part de l'état initial fixe si un départ avait
        été donné, sinon de ce départ avec 2 nd de plus et 2° de gîte en moins.
  \item \textbf{Les commandes du foil :} \py{Foil_Cant}, \py{Foil_Rake} et \py{Foil_Yaw} restent à 0, comme dans
        les cas de navigation vpp36. Il n'y a pas de commande de foil.
\end{itemize}

\begin{subtilite}[dépendance au chemin]
Le solve n'est pas une fonction pure de ses entrées : MDS garde un état interne d'un solve au suivant, par
exemple la fermeture de l'enfoncement. Le 21/10, la même section et le même départ ont donné 14,83 nd
une fois et 17,58 nd après d'autres solves. D'où trois règles :
\begin{itemize}
  \item à chaque mise à jour, les sections sont résolues en partant des états du titulaire
        (départ « chaud ») ;
  \item des départs supplémentaires sont possibles (\py{extra_guesses}), et l'on garde la vitesse convergée
        la plus haute ;
  \item une différence finie compare deux solves \textbf{du même processus} (§\,\ref{sec:sensibilites}).
\end{itemize}
\textbf{Un équilibre lent mais \enquote{réussi}} (26/10). Partie du seul départ fixe, la VPP peut converger
vers un équilibre valide mais lent. Au 4\textsuperscript{e} run, à 48°/16 nd, le meilleur profil a ainsi
donné 15,7 nd (gîte 2,7°, \py{Flat} 0,35) au lieu de 19,2 nd (gîte 5,9°, voiles pleines). L'analyse sur la
distribution (§\,\ref{sec:profil}) part donc de plusieurs départs (\py{profile_analysis.robust_solve}). Dans
la recherche, \textbf{le tout premier solve du titulaire part encore du seul départ fixe} ; les suivants
partent de ses états. Lui donner aussi plusieurs départs est proposé, mais pas fait.
\end{subtilite}

\begin{subtilite}[pas d'extrapolation en attitude]
Les commandes du foil sont nulles, donc $\va$ = (gîte, assiette, dérive). Les bornes de ces trois inconnues,
$[0,8]\times[-1{,}5;1{,}5]\times[-6,6]$, sont exactement la boîte de la RS du foil, et aussi celle du plan
d'expériences du surrogate. Un état VPP ne sort donc jamais du domaine du surrogate en attitude.
\end{subtilite}

\subsection{Le foil dans la VPP}
\label{sec:foilvpp}
Le composant \py{IntensSY_foil_dyn}, dans \py{intenssy_subsystems.py} (dossier
\py{Systems/Ships/Sail/IntensSY} de MDS), fait quatre choses :
\begin{enumerate}
  \item il calcule l'attitude effective :
        \[ a_c = \text{gîte}+\text{Foil\_Cant},\qquad a_r=\text{assiette}+\text{Foil\_Rake},\qquad
           a_y=\text{dérive}+\text{Foil\_Yaw} ; \]
  \item il assemble le vecteur d'entrée de la RS \textbf{par nom}, d'après ses \py{input_names}
        (\py{foil_rs_vector}). La vitesse n'est passée que si la RS déclare l'entrée \py{Speed_kts} ;
  \item il appelle \py{regression(x)}, puis lit \py{prediction} et \py{out('Cpmin')} ;
  \item il redimensionne les 6 premières sorties par $q S=\tfrac12\rho_{\text{eau}} S V^2$, avec
        $S=2{,}275$~m² et $V$ en m/s. Le torseur est appliqué à l'origine du repère du foil, le bord d'attaque
        de la cale basse : $(10{,}5745;\ 3{,}69217;\ -0{,}606084)$ m dans le repère du bateau.
\end{enumerate}

\textbf{Bending} et \textbf{Twisting}, quand la RS les a, sont des sorties de rapport. Elles ne sont pas
multipliées par $qS$ et n'entrent pas dans l'équilibre.

\textbf{Le remplacement par le surrogate.} \py{SurrogateFoilRS} (\py{src/vpp_coupling.py}) a les attributs
que MDS lit sur une \py{VppChebyRS} : \py{input_names}, \py{output_names}, \py{prediction}, \py{eval},
\py{regression}, \py{out}, \py{data}, les drapeaux d'extrapolation. \py{IntensSYVPP.set_foil} le met à la
place de la RS du composant \py{Foil Dyn}. À chaque appel de MDS, avec $x=(a_c,a_r,a_y[,V])$ :
\[
  \text{prediction} = \big(\mu_1..\mu_7\big)(\vz,\va)
  \;+\; \underbrace{r(\vz)\, b_{1..7}(\va)\, V^2}_{\text{correction FSI}},\qquad
  (\text{Bending},\text{Twisting}) = r(\vz)\, b_{8,9}(\va)\, V^2 .
\]
La section $\vz$ est fixée pour un solve.

\begin{nonfait}
\begin{itemize}
  \item \textbf{Pas d'entrée Hvol :} le foil ne voit pas la hauteur d'assise. C'est un choix du 22/10, et la
        RS de production non plus ne la voit pas.
  \item \textbf{La VPP utilise la moyenne du surrogate}, et non sa loi : la variance ne se propage pas dans
        l'équilibre.
  \item \textbf{Cpmin n'entre ni dans l'objectif ni dans une contrainte.} Il est modélisé et rapporté seulement.
\end{itemize}
\end{nonfait}

\subsection{La correction FSI (foil flexible)}
\label{sec:fsi}
La RS FSI de production (\py{CFoil_R9800_S3250_C700_fsi_rs_vpp.json}) vient de bdFoil
\py{CFoil_fsi.ipynb} : AVL × XFOIL avec un modèle de poutre \enquote{supported}, sur 625 points, avec la
vitesse $V\in[7,28]$ nd en 4\textsuperscript{e} entrée. Elle donne l'effet de la déformée pour le profil
actuel.

La déformée suit la charge, donc $V^2$ (structure linéaire), et son effet sur les coefficients aussi :
$C_k(\va,V)=C_k^{\text{rigide}}(\va)+b_k(\va)V^2$. On l'a vérifié sur la RS : l'écart à 7 nd divisé par
$V^2-49$ est constant à 3 \% près jusqu'à 24 nd, et la loi suit la RS à 1,3 \% près au pire sur les forces.
La pente vient d'une régression sur $V_j\in\{7;10{,}5;14;17{,}5;21;24{,}5;28\}$ :
\[
  b_k(\va) = \frac{\sum_j \big(V_j^2-\overline{V^2}\big)\big(C_k^{RS}(\va,V_j)-\overline{C_k^{RS}}\big)}
                  {\sum_j \big(V_j^2-\overline{V^2}\big)^2},
  \qquad
  \delta_k(\vz,\va,V) = r(\vz)\, b_k(\va)\, V^2,\qquad
  r(\vz)=\frac{I_{xx}(\text{profil actuel})}{I_{xx}(\mathcal{D}(\vz))}.
\]
\begin{itemize}
  \item \textbf{Le domaine :} l'attitude est ramenée dans la boîte de la RS. La loi en $V^2$ vaut aussi
        sous 7 nd, sans extrapoler la RS.
  \item \textbf{Le rapport $r$ :} un profil plus raide se déforme moins. La flexion va en $1/EI$, la torsion
        en $1/GJ$, et les deux sont prises proportionnelles à $I_{xx}$ (même famille, même corde, même
        structure). Avec la contrainte $I_{xx}\ge I_{xx}^{\text{actuel}}$, on a $r\le 1$.
  \item \textbf{Le contrôle :} la VPP avec \enquote{partie rigide de la RS + $\delta$} redonne la VPP de
        production à 0,0005 nd près à 45/16, et à 0,007 nd près à 90/14 (test \py{test_fsi_correction_reproduces_the_production_vpp}).
  \item \textbf{L'effet :} la déformée fait gagner 0,075 nd à 45/16 et 0,31 nd à 90/14.
\end{itemize}

\begin{subtilite}[hypothèses de la correction]
Elle suppose deux choses :
\begin{itemize}
  \item à attitude égale, la section porte les charges du profil actuel, à quelques pour cent près ;
  \item la structure est celle de la campagne de production.
\end{itemize}
C'est une correction de premier ordre. Un calcul FSI des profils comparés (bdFoil \py{simulation_fsi}) la
vérifierait.
\end{subtilite}

\subsection{Les sensibilités de vitesse}
\label{sec:sensibilites}
$d_{kc}=\partial V^\star/\partial C_k$ se calcule par différences finies avant, à l'état du titulaire et pour
chaque condition : on décale la sortie $k$ du surrogate de $h_k$, puis on refait le solve depuis l'état de
base :
\[
  d_{kc} \approx \frac{V^\star\big(\mu+h_k\bm{e}_k\big) - V^\star(\mu)}{h_k},\qquad
  h = (5\cdot10^{-4},\,10^{-2},\,10^{-2},\,2\cdot10^{-2},\,3\cdot10^{-3},\,3\cdot10^{-3})
\]
pour $F_x,F_y,F_z,M_x,M_y,M_z$, soit environ 3 \% de chaque coefficient. Le décalage s'applique avant la
correction FSI.

\begin{danslecode}[\py{_vpp_task}, branche \enquote{sensitivity}]
Chaque tâche refait d'abord le solve de base \emph{dans son propre processus}, puis le solve décalé. Avec un
état de base calculé par un autre processus, le premier test court donnait $\partial V/\partial F_x=-32$ au
lieu de $+83$, à cause de la dépendance au chemin.
\begin{itemize}
  \item Une dérivée non calculable (NaN) reprend sa valeur précédente, ou 0 la première fois.
  \item Les sensibilités sont recalculées une mise à jour sur trois.
\end{itemize}
Au 3\textsuperscript{e} run, $d_{F_x}$ allait de 84 (45/12) à 302 (90/14) nd par unité de $F_x$. À la fin
du 4\textsuperscript{e}, il allait de 27 (portant, 135/13) à 170 (près, 48/13). Les autres dérivées sont
plus petites d'un à deux ordres. Une condition sans équilibre garde $d=0$ : elle ne pèse rien dans le
critère.
\end{danslecode}

% ---------------------------------------------------------------------------------------------
\section{Les codes du torseur}
\label{sec:codes}
% ---------------------------------------------------------------------------------------------
\textbf{Le foil} est celui de la RS de production (bdFoil \py{CFoil_generation.ipynb}) :
\begin{itemize}
  \item un arc circulaire \py{cfoil_arc} de rayon 9,8 m et de 3,25 m d'envergure ;
  \item une corde de 0,70 m constante (effilement 1, vrillage nul), 41 stations ;
  \item une tangente d'emplanture à \py{cant_geom} = 20,168°, l'extrémité vers l'intérieur
        (\py{tip_side} $=-1$) ;
  \item une image de paroi en $z=0$ (AVL \py{iZsym} $=+1$ ; NPLLT \py{surface_type} 0).
\end{itemize}

\textbf{Le repère et l'attitude.} Le repère bdFoil a son origine au bord d'attaque de la cale basse, $x$ vers
l'avant, $y$ vers bâbord, $z$ vers le haut. Les moments sont réduits à cette origine, comme dans la RS. Les
entrées passent telles quelles : cant $=a_c$, rake $=a_r$, yaw $=a_y$, sink $=0$.

\textbf{L'écoulement.} $Re=5\cdot10^6$ sur toute l'envergure, $n_{\text{crit}}=1$, transition à 10 \% sur les
deux faces, table d'incidence de $-6°$ à $12°$ par pas de 0,5°.

\textbf{Le profil.} 121 points par face (AVL ne lit pas 161 points et échoue sans rien écrire), avec un bord
de fuite d'épaisseur $0{,}00197\,c$.

\begin{tabularx}{\linewidth}{@{}l>{\raggedright\arraybackslash}X>{\raggedright\arraybackslash}X@{}}
\toprule
Niveau & Modèle & Cpmin \\
\midrule
1 (BF) & ligne portante non plane (Phillips \& Snyder, \py{bdFoil/NonPlanarSolver}), 100 nœuds, rayon de cœur
0,005 corde, polaires NeuralFoil xxxlarge. Temps : 0,2 à 0,4 s. &
20\textsuperscript{e} percentile, sur les points de contrôle, du Cpmin NeuralFoil à l'incidence locale
($1-(u_e/V)^2$, stations de couche limite $x/c\ge1{,}6$ \%) \\
2 (HF) & la même ligne portante avec les polaires XFOIL (\py{npane} 250, \py{niter} 300, 120 s au plus).
Temps : 2 à 13 s pour la polaire d'une nouvelle section, puis 0,1 s par attitude. &
20\textsuperscript{e} percentile du cpmin de la polaire XFOIL à l'incidence locale. La polaire est relue
dans le cache. \\
3 (validation) & AVL (bdFoil core) : traînée induite au plan de Trefftz, traînée de profil = CD de la
polaire XFOIL au cl de chaque bande (modèle \enquote{profile}). Temps : 4,8 s par appel. &
20\textsuperscript{e} percentile des Cpmin de bande (\py{cp_min_perc} 20), comme la RS \\
\bottomrule
\end{tabularx}

\begin{itemize}
  \item \textbf{Les sorties.} Les sorties 1 à 6 sont les coefficients bruts de bdFoil, \py{Cx} … \py{Cmz},
        sans conversion. L'étape 0 a montré qu'ils retrouvent la RS de production : pentes de 0,86 à 0,96 et
        $R^2\ge0{,}996$ pour AVL.
  \item \textbf{Les échecs.} Un calcul raté (NPLLT non convergé, polaire XFOIL ratée, décrochage hors de la
        zone d'emplanture de 2 \% d'envergure) rend des sorties NaN avec son message d'erreur.
  \item \textbf{Les deux niveaux du surrogate.} Sur le profil actuel, NeuralFoil xxxlarge et XFOIL donnent les
        mêmes forces à 0,1 \% près. Ils ne diffèrent vraiment que sur Cpmin : NeuralFoil ne voit pas le pic de
        bord d'attaque.
\end{itemize}

\begin{danslecode}
\py{src/foil_codes.py} : \py{PLANFORM}, \py{FLOW}, \py{LEVELS_XFOIL}, et \py{FoilCodes.evaluate(weights,
attitude, level)}, qui rend \py{(y, info)}. Les backends sont dans \py{pipelines/bdtoolbox_foil/solvers.py}
(\py{NpLltBackend}, \py{AvlCoreBackend}).
\begin{itemize}
  \item Le cache des polaires XFOIL a pour clé le sha256 du fichier \py{.xf} et les réglages. Il publie par
        un seul \py{os.replace}, donc deux processus qui calculent la même section ne se gênent pas : le
        premier qui écrit gagne.
  \item Chaque processus écrit sa propre configuration \py{codes_config_<pid>.json}.
\end{itemize}
\end{danslecode}

% ---------------------------------------------------------------------------------------------
\section{La section : Kulfan, famille admissible, PCA}
\label{sec:section}
% ---------------------------------------------------------------------------------------------
\textbf{Kulfan (CST).} Le poids de bord d'attaque est nul et l'épaisseur de bord de fuite $TE$ fixée. Avec 8
poids par face :
\[
  y_{\text{ext}}(x) = \sqrt{x}(1-x)\sum_{i=0}^{7} u_i B_i(x) + x\,\tfrac{TE}{2},\qquad
  y_{\text{int}}(x) = \sqrt{x}(1-x)\sum_{i=0}^{7} l_i B_i(x) - x\,\tfrac{TE}{2},
\]
où $B_i$ sont les polynômes de Bernstein de degré 7. Aux stations fixées, les coordonnées sont une fonction
affine des 16 poids : $\bm{y}=A\vw+\bm{b}$.

\textbf{Le profil actuel.} Ce sont les poids de référence de l'étude \py{cfoil_intens_kulfan_trim_2}
(\py{CFoilFamily.from_study_config(...).reference_weights()}) : l'épaisseur du MC2 et la cambrure de
référence ($\delta=0{,}108139$), soit 12,25 \% d'épaisseur et $I_{xx}=6{,}7174\cdot10^{-5}$ (corde unité).

\textbf{La famille $\mathcal{F}$} (2\textsuperscript{e} et 3\textsuperscript{e} runs). Une section appartient
à la famille si elle remplit quatre conditions :
\begin{enumerate}
  \item ses 16 poids sont dans une boîte centrée sur le profil actuel, de largeur la moitié de celle de la
        boîte de \py{Dim-Reduction/experiments.ipynb} (\py{KulfanBoxFamily.centered}, \py{notebook_bounds}) ;
  \item ses faces ne se croisent pas : l'épaisseur locale est $\ge0$ sur 2001 stations ;
  \item $I_{xx}\ge I_{xx}^{\text{actuel}}$, avec
        $I_{xx}=\int\big[\tau^3/12+\tau(y_c-y_g)^2\big]dx$, $\tau$ l'épaisseur, $y_c$ la ligne moyenne et $y_g$
        le centre de gravité (\py{section_ixx}) ;
  \item NeuralFoil xxxlarge est confiant : la confiance d'analyse est $\ge0{,}9$ à 0, 2, 4 et 6° (\py{confident}).
\end{enumerate}
Les tirages se font par hypercube latin dans la boîte, puis on rejette les sections hors famille et on retire
jusqu'à en avoir assez (\py{KulfanBoxFamily.sample}).

\textbf{La PCA.}
\begin{itemize}
  \item \textbf{L'ajustement :} sur 4000 sections de la famille. On prend leurs coordonnées à 200 stations
        uniformes, on les centre (sans les standardiser), puis on fait une SVD :
        $Y-\bar{\bm{y}}=USV\tp$.
  \item \textbf{Les scores} sont blanchis : $z_j=(\bm{y}-\bar{\bm{y}})\cdot\bm{v}_j/\sigma_j$.
  \item \textbf{Le décodage} est exact : chaque mode est $\bm{v}_j=A\bm{m}_j$, d'où
        $\mathcal{D}(\vz)=\bar{\vw}+\sum_{j\le k}\sigma_j z_j\bm{m}_j$.
  \item \textbf{Le nombre de modes :} la famille a un rang de 16, et $k=10$ modes portent 99,86 \% de la
        variance géométrique. Ces 10 scores forment l'espace de recherche.
\end{itemize}

\begin{subtilite}[l'admissibilité se vérifie sur la section décodée]
La troncature à 10 modes retire un peu d'$I_{xx}$ : 24 \% des sections projetées sortaient de la famille.
\py{RealFoilProblem.sample_sections} tire donc des sections de la famille, les encode, puis \textbf{ne garde
que les scores dont la section décodée est encore dans la famille}. Celle-ci est la section que voient les
codes. Les bornes des scores (pour normaliser le surrogate) sont les minimums et maximums de 4096 tirages
admissibles.
\end{subtilite}

% ---------------------------------------------------------------------------------------------
\section{Le surrogate du torseur : 7 co-krigeages AR(1)}
\label{sec:surrogate}
% ---------------------------------------------------------------------------------------------
\textbf{La structure.} Il y a un co-krigeage \py{CoKriging} de mfego par sortie : 7 modèles indépendants
($F_x$ … $M_z$, $C_{p,\min}$) sur le \textbf{même plan imbriqué}. Un calcul de code donne en effet tout le
torseur (\py{TorsorCoKriging}, \py{src/vpp_targeted.py}). Les entrées $\vx\in\mathbb{R}^{13}$ sont ramenées
dans $[0,1]^{13}$ avec les bornes du problème.

\textbf{Le modèle auto-régressif} (Kennedy \& O'Hagan), pour chaque sortie :
\[
  Z_1(\vx)=\delta_1(\vx),\qquad Z_2(\vx)=\rho\,Z_1(\vx)+\delta_2(\vx),
\]
où $\delta_1$ et $\delta_2$ sont des GP indépendants, chacun avec une moyenne constante et le noyau
\[
  k(\vx,\vx') = t_1\Big(1+\sqrt5\,r+\tfrac53 r^2\Big)e^{-\sqrt5\,r}+t_2,\qquad
  r^2=\sum_{m=1}^{13}\frac{(x_m-x'_m)^2}{\ell_m^2}
\]
(Matérn 5/2 à longueurs de corrélation par dimension, plus une constante), avec un bruit de pépite.

\textbf{Les hyperparamètres} sont estimés niveau par niveau : sur un plan imbriqué sans bruit, la
vraisemblance se factorise. On les estime par MAP (L-BFGS-B, 3 relances), avec les priors suivants
(sorties normalisées, entrées dans $[0,1]$) :

\begin{tabular}{@{}lll@{}}
\toprule
Paramètre & Prior & Bornes \\
\midrule
$\ell_m$ & InvGamma(3, 2), mode 0,5 & $[10^{-2},10]$ \\
$t_1$ & LogNormale(1, 1) & $[10^{-4},10^2]$ \\
$t_2$ & LogNormale($\ln 10^{-2}$, 2) & $[10^{-8},10^2]$ \\
bruit & LogNormale($\ln 10^{-4}$, 2) & $[10^{-8},10^{-2}]$ \\
$\rho$ & Normale(1, $2^2$) (\enquote{v2}), profilé en forme close, centré & $[-5,5]$ \\
\bottomrule
\end{tabular}

Le niveau 2 s'ajuste sur le résidu $z_2-\rho\,z_1(X_2)$, où $z_1(X_2)$ sont les valeurs \emph{observées} du
niveau 1 aux points HF. Le résidu est centré, donc sa moyenne est nulle quel que soit $\rho$. Pour $\Theta$
fixé, $\hat\rho$ est l'optimum en forme close sous le prior normal.

\textbf{La prédiction} conditionne $Z_2(\vx)$ sur \emph{toutes} les observations $\bm{z}$ (des deux
niveaux), avec la matrice de covariance jointe $V$, factorisée une fois par Cholesky ($V=LL\tp$) :
\[
  \mu(\vx)=m_2+\bm{k}(\vx)\tp V^{-1}(\bm{z}-\bm{m}),\qquad
  s^2(\vx)=c_{22}(\vx,\vx)-\bm{k}(\vx)\tp V^{-1}\bm{k}(\vx),
\]
avec $\Cov(Z_i(\vx),Z_j(\vx'))=\sum_{m\le\min(i,j)}P(m,i)P(m,j)\,\sigma_{y,m}^2\,k_m(\vx,\vx')$ et
$P(1,2)=\rho$.

\begin{danslecode}[ce qui coûte, et comment c'est réparti]
\begin{itemize}
  \item \py{fit(refit=True)} : de nouveaux hyperparamètres (MAP). \py{fit(refit=False)} : on garde les
        hyperparamètres, et l'on reconditionne seulement (Cholesky de $V$, $\alpha=V^{-1}(\bm z-\bm m)$).
        Les 7 modèles tournent dans 7 fils (\py{TorsorCoKriging.map_models}), et le résultat est identique
        au calcul en série.
  \item \textbf{Le rythme des réajustements :} de nouveaux hyperparamètres dès que le plan a grossi de 25 \%
        depuis le dernier ajustement (\py{refit_growth} = 1,25), et plus du tout au-delà de
        \py{--max-refit-size} points (2\,500 au 4\textsuperscript{e} run). Un ajustement coûte $O(n^3)$ :
        environ 2,5 min à 1\,045 points au 3\textsuperscript{e} run, de l'ordre de l'heure vers 3\,000. Au 3\textsuperscript{e} run, cela s'est produit à 314, 408, 518, 651, 830 et 1045 points ;
        au 4\textsuperscript{e}, de 314 à 2\,098 points, puis les hyperparamètres sont restés figés. Sinon, on
        reconditionne seulement.
  \item \textbf{Pour la VPP, la moyenne seule :} \py{TorsorCoKriging.predict} appelle \py{latent_mean}, soit
        $m_2+\bm{k}(\vx)\tp\alpha$, en $O(n)$. La variance demanderait une résolution triangulaire en
        $O(n^2)$ par point, et la VPP appelle le foil des centaines de fois par solve.
  \item \textbf{Les workers VPP} reçoivent une copie sans le facteur de Cholesky (\py{mean_only_copy}), qui
        ne sait prédire que les moyennes.
\end{itemize}
\end{danslecode}

% ---------------------------------------------------------------------------------------------
\section{Le critère de raffinement ciblé par la VPP}
\label{sec:critere}
% ---------------------------------------------------------------------------------------------
\textbf{L'idée.} À un état d'équilibre, une erreur $\Delta C_k$ sur le coefficient $k$ change la vitesse de
$d_k\Delta C_k$ environ. Les modèles sont indépendants, donc au premier ordre
\[
  \Var V \approx \sum_{k=1}^{6} d_k^2\, s_k^2(\vx).
\]
On cherche à réduire cette variance \textbf{en moyenne sur les points d'intégration}, qui entourent les
états VPP des sections en compétition. Cpmin a un poids nul : il n'agit pas sur la vitesse.

\textbf{Les points d'intégration $\{\vx^{\text{int}}_i\}_{i=1}^N$} (\py{RealFoilProblem.cloud}). À chaque
mise à jour, on prend les états des 4 sections aux conditions résolues : 20 au 3\textsuperscript{e} run
(5 conditions), jusqu'à 72 au 4\textsuperscript{e} (18 nœuds). Autour de chaque état $s=(\vz_s,\va_s)$, on
tire un nuage de Sobol de 16 points :
\[
  \vx=\Big(\vz_s\ \big[+R\,\bm g/\sqrt{k}\ \text{un point sur deux}\big],\;
  \operatorname{clip}\big(\va_s+\bm\xi\odot(1°;\,0{,}3°;\,1°)\big)\Big),\qquad
  \bm\xi\in[-1,1]^3,\ \bm g\sim\mathcal N(0,I).
\]
Sans recherche locale, la section reste fixée. Avec elle (\py{--local-fraction} $>0$), un point sur deux
déplace aussi la section d'une distance typique $R$, le rayon de la région de confiance
(§\,\ref{sec:optimiseur}), et l'on retire les points hors de la famille. Au 3\textsuperscript{e} run,
$N=320$ ; au 4\textsuperscript{e}, $N\approx1\,150$.

\textbf{Le poids de chaque point} (\py{point_weights}, depuis le 23/10). Avec $c(i)$ la condition de l'état du
point $i$ et $n_{s(i)}$ le nombre de points gardés pour cet état :
\[
  \omega_{ik}=\frac{w_{c(i)}\;d_{k\,c(i)}^2}{n_{s(i)}} .
\]
Chaque état compte pour la probabilité de sa condition, quel que soit le nombre de ses points. Les
sensibilités sont celles du titulaire à cette condition, appliquées aussi aux états des concurrentes.
Jusqu'au 3\textsuperscript{e} run, le poids était le même pour toutes les conditions :
$w_k=\frac15\sum_c d_{kc}^2$.

\textbf{Le gain exact d'une observation} (ALC/IMSE, \py{batched_imse_gains}). Soit un candidat $\vx_m$ et un
niveau $l$. Le plan doit rester imbriqué : on ajoute les observations $(\vx_m,j)$ pour tous les niveaux
$j\le l$ qui manquent en $\vx_m$. Avec $\bm{c}_i$ la covariance a posteriori entre $Z_2(\vx^{\text{int}}_i)$
et ces observations, et $S$ leur covariance a posteriori (bruit compris), la réduction exacte de la
variance de $Z_2$ en $\vx^{\text{int}}_i$ est $\bm{c}_i\tp S^{-1}\bm{c}_i$. D'où, pour le co-krigeage $k$,
\[
  G_k(\vx_m,l)=\sum_{i=1}^{N}\omega_{ik}\,\bm{c}_i\tp S^{-1}\bm{c}_i ,
  \qquad G(\vx_m,l)=\sum_{k\le6}G_k(\vx_m,l)\quad[\text{nd}^2].
\]
Le calcul se fait pour tous les candidats d'un coup (tableaux à un axe de plus), avec les blocs
$W=L^{-1}C(\D,\cdot)$.

\textbf{C'est un critère ALC, pas ALM.} L'ALM (MacKay) calcule là où la variance prédite est maximale :
c'est la règle de Le Gratiet \& Cannamela pour le point. L'ALC (Cohn) calcule là où l'observation fait le
plus baisser la variance sur des points de référence, ici les points d'intégration pondérés. Le code a les
deux (\py{SequentialCoKriging(point_criterion="variance"|"alc")}), et la boucle VPP n'utilise que l'ALC.

\textbf{Le rendement et le choix.}
\[
  (\vx^\star,l^\star)=\argmax_{m,\,l}\;\frac{G(\vx_m,l)}{C(\vx_m,l)},\qquad
  C(\vx_m,l)=\sum_{j\le l,\ j\text{ manquant}} c_j,\qquad c=(1,2).
\]
Pour un point nouveau, le niveau 2 coûte $1+2=3$. Pour un point déjà calculé en BF, il coûte 2 : depuis le
23/10, cela n'arrive qu'à l'intérieur d'un lot, quand un candidat pris en BF est repris en HF.

\begin{subtilite}[le rapport de coût 1 : 2]
Les codes d'un lot tournent en même temps. Un choix coûte surtout le temps de sélection, de 1 à 6 s, à peu
près le même aux deux niveaux. Un calcul HF n'allonge le lot que si sa section est nouvelle (la polaire
XFOIL). Au 2\textsuperscript{e} run, en séquentiel avec AVL, le rapport valait 1 : 6.
\end{subtilite}

\textbf{Les candidats} (\py{VPPSectionOptimizer.candidates}).
\begin{itemize}
  \item D'autres nuages autour des états, tirés comme les points d'intégration avec une autre graine :
        $\lfloor512/N_{\text{états}}\rfloor$ par état, soit 25 au 3\textsuperscript{e} run et 7 au
        4\textsuperscript{e}.
  \item On retire les candidats à moins de $10^{-3}$ (cube unité, norme infinie) d'un calcul raté.
  \item Jusqu'au 3\textsuperscript{e} run, il y avait aussi 256 points de tout le domaine et au plus 512 points
        BF sans HF. Ils ont été retirés le 23/10 : au 3\textsuperscript{e} run, 100 \% des choix étaient près
        des états, et 2 HF sur 252 seulement sur un ancien point BF.
\end{itemize}

% ---------------------------------------------------------------------------------------------
\section{Les lots en parallèle : le \enquote{kriging believer}}
\label{sec:lots}
% ---------------------------------------------------------------------------------------------
On veut 16 choix par lot pour faire tourner les codes en parallèle. Ce n'est pas un choix des 16 meilleurs
rendements d'un coup : ces points seraient presque tous voisins et redondants. On choisit plutôt l'un après
l'autre. Après chaque choix, on conditionne le co-krigeage sur le point choisi comme s'il avait été
calculé, avec pour valeur provisoire la \textbf{moyenne prédite du niveau concerné}.

\textbf{Pourquoi c'est exact pour notre critère.} La covariance a posteriori d'un GP,
\[
  k_n(\bm{a},\bm{b})=k(\bm{a},\bm{b})-\bm{k}(\bm{a},X)K^{-1}\bm{k}(X,\bm{b}),
\]
ne dépend que des \emph{positions} des observations, et non de leurs valeurs. Le gain $G_k$ du choix suivant voit donc exactement
la réduction de variance apportée par les choix précédents. La valeur provisoire laisse les moyennes
inchangées (test \py{test_kriging_believer_placeholder}). Le lot s'étale tout seul.

\begin{algorithm}[h]
\caption{Un lot (\py{VPPSectionOptimizer.step})}
\begin{algorithmic}[1]
\State $n\gets$ taille du plan ; \textbf{si} $n\ge1{,}25\,n_{\text{dernier ajustement}}$ et $n\le n_{\max}$ \textbf{alors} nouveaux hyperparamètres \textbf{sinon} reconditionnement
\State mise à jour de l'optimiseur (§\,\ref{sec:optimiseur}) : titulaire, états, points d'intégration, $d_{kc}$
\State poids $\omega_{ik}\gets w_{c(i)}d_{k\,c(i)}^2/n_{s(i)}$ ; candidats $\{\vx_m\}$ et niveaux manquants
\For{$b=1..16$}
  \State $(\vx_b,l_b)\gets$ meilleur rendement (§\,\ref{sec:critere}) ; marquer ses niveaux $\le l_b$ comme faits
  \If{$b<16$}
     \For{chaque niveau $j\le l_b$ manquant en $\vx_b$} ajouter $(\vx_b,j)$ avec la valeur $\mu^{(j)}(\vx_b)$ \EndFor
     \State reconditionner les 7 modèles (hyperparamètres fixes)
  \EndIf
\EndFor
\State retirer les valeurs provisoires (\py{truncate})
\State lancer en parallèle tous les calculs (point, niveau) du lot (28 processus)
\For{chaque choix $b$}
   \State ajouter ses résultats, niveau 1 d'abord. Un résultat NaN, ou un HF dont le BF manque, est rejeté, et le point rejoint la liste des échecs
\EndFor
\end{algorithmic}
\end{algorithm}

% ---------------------------------------------------------------------------------------------
\section{L'optimiseur de section}
\label{sec:optimiseur}
% ---------------------------------------------------------------------------------------------
Il y a une mise à jour avant chaque lot (\py{update_every} = 1, \py{VPPSectionOptimizer.update}).
\begin{enumerate}
  \item \textbf{La publication} du surrogate courant pour les workers VPP.
  \item \textbf{Les sensibilités} $d_{kc}$, une mise à jour sur trois.
  \item \textbf{Le criblage} (\py{screen}). On tire 256 sections admissibles nouvelles, et on y ajoute le
        titulaire $\vz_{\text{tit}}$ :
        \begin{itemize}
          \item une part \py{local_fraction} (0,5 au 4\textsuperscript{e} run, soit 128 sections) dans la
                région de confiance, $\vz_{\text{tit}}+R\,\bm g/\sqrt{k}$ avec $\bm g\sim\mathcal N(0,I)$,
                gardées si la section décodée est dans la famille (\py{local_sections}) ;
          \item les autres dans toute la famille, avec la graine $\text{screen\_seed}+31\,n_{\text{mises à
                jour}}$. \py{screen_seed} est propre au run (tirée au hasard par défaut, \py{--screen-seed}) :
                avant le 23/10, elle valait la graine du run, et les runs 2 et 3 ont criblé les mêmes sections.
        \end{itemize}
        On garde fixe l'attitude $\va_c$ des états résolus du titulaire et sa vitesse $V_c$. Le changement de
        vitesse au premier ordre vaut alors
        \[
          \Delta V_c(\vz)=\sum_{k\le6} d_{kc}\big[\mu_k(\vz,\va_c)-\mu_k(\vz_{\text{tit}},\va_c)\big]
          + \big(r(\vz)-r(\vz_{\text{tit}})\big)\sum_{k\le6} d_{kc}\, b_k(\va_c)V_c^2,
        \]
        où le second terme est la différence des corrections FSI. On pondère par la probabilité des
        conditions :
        \[
          \mu(\vz)=\frac{\sum_c w_c\,\Delta V_c(\vz)}{\sum_c w_c},\qquad
          \sigma(\vz)=\frac{\big(\sum_c w_c^2\sum_{k\le6} d_{kc}^2\,s_k^2(\vz,\va_c)\big)^{1/2}}{\sum_c w_c},\qquad
          \text{score}(\vz)=\mu(\vz)+\kappa\,\sigma(\vz),\quad\kappa=1 .
        \]
        Le terme $\kappa\sigma$ est un optimisme face à l'incertitude : une section dont la vitesse est
        incertaine mérite d'être regardée. On retient les 3 meilleurs scores.
  \item \textbf{La VPP complète} sur ces 3 sections et sur le titulaire, à toutes les conditions (20 solves
        au 3\textsuperscript{e} run, 72 au 4\textsuperscript{e}), en parallèle, avec départ aux états du
        titulaire. La section de meilleure vitesse moyenne pondérée $J$ (§\,2) devient le titulaire.
  \item \textbf{La région de confiance} (recherche locale) :
        \[
          R\leftarrow\min(R_{\max},\,1{,}5\,R)\ \text{si le titulaire change},\qquad
          R\leftarrow\max(R_{\min},\,0{,}7\,R)\ \text{sinon},\qquad [R_{\min},R_{\max}]=[0{,}1;\,3],
        \]
        en partant de $R=1$ (\py{--radius}, en écarts types des scores blanchis). C'est la règle des
        méthodes de région de confiance (Conn, Gould \& Toint, 2000) et de TuRBO (Eriksson et al., 2019).
  \item \textbf{Les points d'intégration} : nuages de Sobol autour des états des 4 sections, dont la moitié
        déplace aussi la section du rayon $R$ (§\,\ref{sec:critere}).
\end{enumerate}

\begin{subtilite}[$w_c$ dans le critère, $w_c^2$ dans le criblage]
Le critère réduit $\sum_c w_c\Var V_c$, une moyenne des variances : chaque condition probable doit être
connue. Le $\sigma$ du criblage est l'écart type de la \emph{moyenne} pondérée, en supposant les erreurs
indépendantes d'une condition à l'autre. Elles sont en fait corrélées (un même défaut de traînée vaut à
toutes les conditions), et ce $\sigma$ est trop petit : l'exploration est un peu plus timide que $\kappa=1$
ne le laisse croire.
\end{subtilite}

\textbf{Le départ et la fin.}
\begin{itemize}
  \item \textbf{Le départ :} le premier titulaire est le profil actuel projeté, et ses états sont résolus à la
        première mise à jour.
  \item \textbf{La boucle} s'arrête quand le temps écoulé depuis le début du run dépasse la limite. Le lot en
        cours se termine.
  \item \textbf{La fin :} un dernier ajustement des hyperparamètres (sauf au-delà de \py{--max-refit-size}),
        puis une dernière mise à jour. Celle-ci fixe le titulaire final avec le modèle final.
\end{itemize}

\begin{subtilite}[ce que la recherche explore vraiment]
Le titulaire ne bouge que vers une section \emph{tirée} au criblage : la recherche trie des sections
aléatoires, environ 13\,000 en 52 mises à jour au 3\textsuperscript{e} run. Ce n'est pas un optimum local au
sens strict.
\begin{itemize}
  \item \textbf{Jusqu'au 3\textsuperscript{e} run}, les tirages dépendaient de la seule graine du run. Les
        2\textsuperscript{e} et 3\textsuperscript{e} runs, de même graine, même famille et même PCA, ont donc
        criblé les mêmes sections. Ils ont choisi le même profil à la même mise à jour (la
        14\textsuperscript{e}), alors que leurs modèles différaient : AVL rigide d'un côté, NPLLT × XFOIL
        flexible de l'autre.
  \item \textbf{Au 4\textsuperscript{e} run}, avec la région de confiance, les sections locales ont pris les 3
        places de concurrentes dans 116 mises à jour sur 121, dès la deuxième. Le
        rayon est monté à 3 tant que le titulaire progressait, puis il est descendu jusqu'à 0,1, atteint à
        223 min. Il y a eu 48 changements de titulaire ; après les 30 premières mises à jour (74 min), aucun
        n'a apporté plus de 0,006 nd.
\end{itemize}
\end{subtilite}

% ---------------------------------------------------------------------------------------------
\section{Le plan de départ}
% ---------------------------------------------------------------------------------------------
\begin{itemize}
  \item \textbf{Le niveau 1 :} $n_1=256$ points, avec des sections admissibles tirées (§\,\ref{sec:section})
        et une attitude par hypercube latin dans $[0,8]\times[-1{,}5;1{,}5]\times[-6,6]$.
  \item \textbf{Le niveau 2 :} $n_2=64$ points choisis parmi les points BF réussis, par maximin glouton dans
        le cube unité (\py{greedy_maximin}). Le premier est le plus proche du centre, puis on ajoute à chaque
        fois le point le plus loin des points déjà pris. Le plan est donc imbriqué.
  \item \textbf{Les échecs} sont écartés. Au 3\textsuperscript{e} run, il y a eu 253/256 BF et 61/64 HF
        réussis, en 1,4 min de calcul parallèle. Le 4\textsuperscript{e} run a le même plan de départ.
\end{itemize}

% ---------------------------------------------------------------------------------------------
\section{L'architecture logicielle et parallèle}
\label{sec:architecture}
% ---------------------------------------------------------------------------------------------
\begin{tabularx}{\linewidth}{@{}l>{\raggedright\arraybackslash}X@{}}
\toprule
Fichier & Rôle \\
\midrule
\py{scripts/run_real_vpp.py} & orchestration en 5 étapes : famille et PCA, plan de départ, boucle,
validation, résultats et figures \\
\py{src/section_pca.py} & Kulfan, \py{KulfanBoxFamily}, \py{SectionPCA}, \py{section_ixx} \\
\py{src/meta_foil.py} & lecteur de RS de Chebyshev (\py{ChebyshevRS}), filtre \py{confident} \\
\py{src/foil_codes.py} & \py{FoilCodes} : les 3 niveaux, Cpmin \\
\py{pipelines/bdtoolbox_foil/} & backends NPLLT, AVL et XFOIL (bdFoil) \\
\py{src/vpp_coupling.py} & \py{FSICorrection}, \py{SurrogateFoilRS}, \py{IntensSYVPP} (solve, sensibilités) \\
\py{mfego/src/cokriging.py} & \py{CoKriging} AR(1) \\
\py{src/vpp_targeted.py} & \py{TorsorCoKriging}, \py{latent_mean}, \py{choose} \\
\py{src/sequential_design.py} & \py{batched_imse_gains} (gain ALC exact) \\
\py{src/vpp_optimize.py} & \py{RealFoilProblem}, \py{CodeRunner}, \py{VPPPool}, \py{VPPSectionOptimizer}, \py{validate} \\
\py{src/run_dir.py}, \py{scripts/real_run_figures.py} & dossier de run, figures \\
\py{src/operating_profile.py}, \py{profiles/*.json} & distribution d'usage : mélange de modes, nœuds et poids \\
\py{scripts/profile_analysis.py} & comportement des sections sur une distribution d'usage, solves à départs multiples (§\,\ref{sec:profil}) \\
\bottomrule
\end{tabularx}

\textbf{Les processus} (fig.~\ref{fig:architecture}), sur un Ryzen 9 5950X de 16 cœurs (32 logiques) avec
128 Go de mémoire.
\begin{itemize}
  \item \textbf{Le processus principal :} la boucle, le surrogate et le critère. Il a 2 fils BLAS, et les
        7 modèles tournent en 7 fils pendant l'ajustement, le conditionnement, les variances et les gains.
  \item \textbf{\py{CodeRunner} : 28 processus.} Chacun a son \py{FoilCodes}. Le pool est créé au premier lot
        de plusieurs points, puis gardé pendant tout le run. Un point seul tourne dans le processus principal.
  \item \textbf{\py{VPPPool} : 30 processus.} Chacun a son \py{IntensSYVPP}, avec la correction FSI. À chaque
        mise à jour, le processus principal écrit \py{data/surrogate_v<n>.pkl}, la copie \enquote{moyennes
        seules}. Un worker ne relit le fichier que si la version a changé, et le fichier précédent est
        supprimé.
  \item \textbf{Les workers des deux pools} démarrent avec un seul fil BLAS (variables d'environnement posées
        avant leur création). Les deux pools ne travaillent jamais en même temps.
\end{itemize}

\textbf{Le dossier d'un run}, \py{runs/<date>_<nom>/} :
\begin{itemize}
  \item \py{config.json} : tous les réglages ;
  \item \py{run.log} ;
  \item \py{data/evaluations.csv} : un calcul de code par ligne, avec le niveau, les entrées, les sorties,
        le temps et l'erreur ; aussi \py{data/pca.json} et \py{data/codes/} (configurations, cache des
        polaires) ;
  \item \py{results/} : \py{history.csv} (un choix par ligne), \py{updates.json}, \py{validation.csv},
        \py{best_section.json}, \py{summary.txt}, \py{surrogate.pkl} (le surrogate final, copie
        \enquote{moyennes seules}), et \py{profile_<nom>.csv|json} si une distribution d'usage est
        analysée ;
  \item \py{figures/} : convergence, points, sections, validation, vitesses, et \py{profile_<nom>}.
\end{itemize}

% ---------------------------------------------------------------------------------------------
\section{La validation finale}
% ---------------------------------------------------------------------------------------------
Pour chaque section validée (le meilleur profil, le profil de départ et les meilleurs profils des runs
antérieurs, donnés par \py{--compare}) :
\begin{enumerate}
  \item on résout la VPP aux conditions du run avec le surrogate final et la FSI ;
  \item on lance les codes de niveaux 1, 2 et 3 à ces états, en parallèle ;
  \item on corrige la vitesse au premier ordre de l'écart mesuré entre le surrogate et le code :
        \[
          V_{\text{code}}\approx V_{\text{surrogate}}-\sum_{k\le6} d_{kc}\big(\mu_k-C_k^{\text{code}}\big).
        \]
        Le code est soit le niveau 2, soit AVL. On compare aussi l'erreur à l'écart type annoncé,
        $\big(\sum_k d_{kc}^2 s_k^2\big)^{1/2}$ ;
  \item on compare les sections par leur vitesse corrigée, en moyenne pondérée par $w_c$ sur les conditions
        résolues pour \emph{toutes} les sections validées (\py{run_real_vpp.summarize}, corrigé le 26/10 :
        avant, le résumé comparait la vitesse de départ vue par le premier surrogate à celle du titulaire
        final).
\end{enumerate}
On n'a pas résolu la VPP avec un code dans la boucle : avec AVL, cela prendrait environ 30 min par condition.

\textbf{Le résultat du 3\textsuperscript{e} run.}
\begin{itemize}
  \item \textbf{Le gain} du meilleur profil sur le profil actuel vaut +0,241 nd en moyenne vu par
        NPLLT × XFOIL, et +0,217 nd vu par AVL × XFOIL. Il tient aux 5 conditions, de +0,02 à +0,40 nd.
  \item \textbf{Le Cpmin} passe de $-2{,}57$ à $-0{,}87$.
  \item \textbf{La précision :} l'erreur du surrogate contre NPLLT × XFOIL vaut 0,010 nd RMS au meilleur
        profil.
\end{itemize}

\textbf{Le résultat du 4\textsuperscript{e} run} (18 nœuds de \py{guess_bimodal}, dont 13 résolus pour les
trois sections validées, soit 72 \% de la probabilité).
\begin{itemize}
  \item \textbf{Le profil :} encore moins cambré que celui des runs 2 et 3 (poids de cambrure moyen 0,081,
        contre 0,090 et 0,108 pour le profil actuel), de 12,54 \% d'épaisseur, avec $I_{xx}$
        \textbf{exactement à la contrainte}. L'optimum est sur cette frontière : le niveau de la contrainte
        est un choix de conception qui fixe le résultat.
  \item \textbf{Le gain} sur le profil actuel vaut +0,161 nd vu par NPLLT × XFOIL et +0,156 nd vu par AVL. Il
        vaut +0,056 nd sur le meilleur profil des runs 2 et 3. Il est positif aux 13 conditions : de +0,01 à
        +0,03 nd au près par petit temps, de +0,17 à +0,42 nd au portant et à 16 nd au près.
  \item \textbf{La précision :} 0,0001 nd RMS contre NPLLT × XFOIL au meilleur profil, très raffiné, et
        0,008 nd au profil actuel ; 0,018 à 0,022 nd contre AVL. L'écart type annoncé est un peu optimiste :
        $|z|\le2$ sur 78 à 83 \% des états, au lieu de 95 \%.
\end{itemize}

% ---------------------------------------------------------------------------------------------
\section{L'analyse sur une distribution d'usage}
\label{sec:profil}
% ---------------------------------------------------------------------------------------------
\textbf{Le principe.} Le bateau navigue selon une distribution de (TWA, TWS), en général bimodale : un mode
au près, un mode au portant. On ne la connaît pas encore. Elle se donne donc \textbf{à chaque run}, dans un
fichier JSON, et elle sert deux fois :
\begin{itemize}
  \item \textbf{pour analyser} les sections retenues (le meilleur profil, le profil de départ et les profils
        comparés), avec \py{--profile} ou \py{scripts/profile_analysis.py} (décision du 22/10) ;
  \item \textbf{pour chercher}, avec \py{--search-profile} (23/10) : ses nœuds deviennent les conditions du
        run, et leurs poids $w_c$ pondèrent l'objectif (§\,2), le criblage (§\,\ref{sec:optimiseur}) et le
        critère (§\,\ref{sec:critere}). Le co-krigeage, lui, ne change pas.
\end{itemize}

\textbf{La distribution} (\py{src/operating_profile.py}) est un mélange de modes. Chaque mode $m$ a un poids
$\pi_m$, une loi normale tronquée de TWA et une loi normale tronquée de TWS, indépendantes dans le mode :
\[
  p(\mathrm{TWA},\mathrm{TWS})=\sum_m \pi_m\,\mathcal{N}_{[a_m,b_m]}(\mathrm{TWA};\bar\theta_m,\sigma_{\theta,m})\,
  \mathcal{N}_{[c_m,d_m]}(\mathrm{TWS};\bar v_m,\sigma_{v,m}).
\]
\begin{itemize}
  \item \textbf{La quadrature.} Elle est produit par mode : $n$ nœuds par variable, aux points de probabilité
        médiane de $n$ strates équiprobables (les quantiles $(i+\frac12)/n$ de la loi tronquée), chacun de
        poids $\pi_m/n^2$. On a donc $\mathbb{E}_p[Q]\approx\sum_{\text{nœuds}} w\,Q$.
  \item \textbf{Les bornes.} Les nœuds restent dans les bornes : pas de queue hors du domaine de la VPP.
  \item \textbf{Des nœuds explicites.} Un fichier peut aussi donner ses nœuds directement, par exemple un
        histogramme de conditions enregistrées. Les poids sont alors normalisés à 1.
\end{itemize}

\textbf{L'hypothèse actuelle} (\py{profiles/guess_bimodal.json}, à remplacer) :
\begin{itemize}
  \item un mode \enquote{près}, TWA $45\pm4°$ sur $[38,52]$ ;
  \item un mode \enquote{portant}, TWA $135\pm7°$ sur $[120,148]$ ;
  \item le même poids pour les deux, TWS $13\pm3$ nd sur $[8,18]$ dans les deux modes ;
  \item $3\times3$ nœuds par mode, soit 18 nœuds.
\end{itemize}

\textbf{L'analyse} (\py{scripts/profile_analysis.py}, ou l'option \py{--profile} de \py{run_real_vpp.py}) se
fait en quatre étapes :
\begin{enumerate}
  \item on recharge le surrogate final du run ; un reconditionnement, sans changer les hyperparamètres, lui
        rend ses variances ;
  \item on résout la VPP (avec la FSI si le run l'avait) à chaque nœud, pour chaque section, dans le pool
        VPP, \textbf{avec des départs multiples} (\py{robust_solve}, 26/10). Une première passe part du
        départ fixe et de deux départs génériques, l'un au près toute la toile dehors, l'autre au portant.
        Une seconde passe part des équilibres trouvés par toutes les sections au même nœud : des sections
        voisines ont des équilibres voisins. On garde le solve réussi le plus rapide ;
  \item on calcule à chaque nœud les sensibilités, puis l'écart type de la vitesse prédite
        $\big(\sum_k d_k^2 s_k^2\big)^{1/2}$. Le surrogate a été raffiné autour des états des conditions du
        run, et un nœud loin d'eux (au portant) est moins sûr ;
  \item on prend les espérances \textbf{sur les nœuds résolus par toutes les sections}, pour comparer les
        mêmes conditions : $\mathbb{E}_p[V]$, la vitesse équivalente en temps $1/\mathbb{E}_p[1/V]$, l'écart
        type, les mêmes par mode, et le gain sur le profil de départ. La probabilité résolue par chaque
        section est donnée à part.
\end{enumerate}

\begin{subtilite}[des nœuds sans équilibre]
La VPP de l'Intens SY n'a pas d'équilibre au portant profond par petit temps. La gîte bute sur sa borne de
0°, avec un résidu non nul, et c'est vrai aussi avec la RS de production. Elle n'en a pas non plus par vent
fort (45/20, 90/20 : vitesse à sa borne de 24 nd ou \py{Flat} au minimum). Ces nœuds sont signalés, et non
cachés. Ils relèvent du réglage de la VPP (bornes, configuration de foil au portant), pas du surrogate.
\end{subtilite}

\textbf{Les résultats} (refaits le 26/10 avec les départs multiples).
\begin{itemize}
  \item \textbf{Le 3\textsuperscript{e} run}, analysé après coup sur le surrogate
        \py{data/surrogate_v52.pkl}. Les deux profils ensemble résolvent 72 \% de la probabilité, tout le
        près et 22 \% sur les 50 \% du portant. Le gain du meilleur profil sur le profil actuel vaut
        \textbf{+0,093 nd} : +0,076 nd au près et +0,13 nd au portant. Sur les 5 conditions du run, centrées
        sur le reaching, il valait +0,24 nd. L'écart type du surrogate est d'environ 0,006 nd, une estimation
        du modèle seul : ce run n'a pas validé le portant contre les codes.
  \item \textbf{Le 4\textsuperscript{e} run}, dont la recherche visait cette distribution. Le gain vaut
        \textbf{+0,152 nd} : +0,103 nd au près, +0,262 nd au portant, positif aux 13 nœuds résolus (de +0,006
        à +0,42 nd). Le meilleur profil des runs 2 et 3 y gagne +0,100 nd.
  \item \textbf{Une erreur corrigée.} Avec le seul départ fixe, l'analyse du 4\textsuperscript{e} run donnait
        $-0{,}23$ nd, contre +0,16 nd à la validation : à 48/16, le meilleur profil tombait sur l'équilibre
        lent du §\,\ref{sec:vpp}.
\end{itemize}

% ---------------------------------------------------------------------------------------------
\section{Les réglages des 3\textsuperscript{e} et 4\textsuperscript{e} runs}
% ---------------------------------------------------------------------------------------------
\begin{tabularx}{\linewidth}{@{}lX@{}}
\toprule
Réglage & Valeur \\
\midrule
commande, 3\textsuperscript{e} run & \py{run_real_vpp.py --box current --box-scale 0.5 --hf npllt_xfoil --hf-cost 2 --fsi --batch 16 --workers 28 --vpp-workers 30 --n-low 256 --n-high 64 --time-limit-min 65 --compare <best_section.json du 2e run>} \\
commande, 4\textsuperscript{e} run & la même, avec \py{--time-limit-min 480 --search-profile profiles/guess_bimodal.json --local-fraction 0.5 --radius 1.0 --max-refit-size 2500 --profile profiles/guess_bimodal.json} et \py{--compare <best_section.json du 3e run>} (lancée par \py{runs/2026-10-23_weekend/orchestrate.ps1}) \\
conditions & 3\textsuperscript{e} : 5 conditions, $w_c=\frac15$ ; 4\textsuperscript{e} : 18 nœuds de \py{guess_bimodal}, $w_c=\frac1{18}$ \\
PCA & $k=10$, ajustée sur 4000 sections, graine 0 \\
plan de départ & 256 BF, 64 HF \\
lots & 16 choix, rapport de coût 1 : 2 \\
mise à jour & à chaque lot ; sensibilités une fois sur trois ; criblage de 256 sections, 3 retenues, $\kappa=1$ ; 4\textsuperscript{e} : la moitié dans la région de confiance, \py{screen_seed} 875768254 \\
points d'intégration & 16 par état, demi-largeurs $(1°;\,0{,}3°;\,1°)$ ; 4\textsuperscript{e} : un point sur deux déplace la section du rayon $R$ \\
candidats & 3\textsuperscript{e} : 25 par état (500), 256 dans tout le domaine, $\le512$ BF sans HF ; 4\textsuperscript{e} : 7 par état, nuages seulement \\
réajustement & croissance de 25 \% du plan ; 4\textsuperscript{e} : plus au-delà de 2\,500 points \\
co-krigeage & Matérn 5/2 à longueurs par dimension + constante, MAP, $\rho\sim\mathcal{N}(1,2^2)$ sur $[-5,5]$, 3 relances \\
FSI & RS \py{CFoil_R9800_S3250_C700_fsi_rs_vpp.json}, rapport $I_{xx}$ \\
arrêt & 3\textsuperscript{e} : 65 min (fin réelle 69,7 min) ; 4\textsuperscript{e} : 480 min (boucle de 489 min, 120 lots, 1\,920 choix) \\
\bottomrule
\end{tabularx}

% ---------------------------------------------------------------------------------------------
\section{Limites connues et évolutions prévues}
% ---------------------------------------------------------------------------------------------
\begin{enumerate}
  \item \textbf{La recherche locale reste par tirages} (fait le 23/10 : région de confiance, graine de
        criblage propre). Elle tire des sections au hasard dans le voisinage, sans gradient. Au
        4\textsuperscript{e} run, cela a suffi : la titulaire ne gagnait plus rien après 75 min. Un gradient
        sur le surrogate irait plus vite, sans doute sans aller plus loin.
  \item \textbf{L'optimum est sur la contrainte} $I_{xx}\ge I_{xx}^{\text{actuel}}$ au 4\textsuperscript{e}
        run. Son niveau est un choix de conception, à arrêter avant un prochain run.
  \item \textbf{La distribution d'usage est une hypothèse} (\py{guess_bimodal}, cherchée depuis le 23/10), et
        28 \% de sa probabilité n'a pas d'équilibre VPP (portant profond par petit temps). Il faudra la
        vraie distribution, et une VPP résoluble sur tout son domaine.
  \item \textbf{Le premier solve du titulaire} part du seul départ fixe de la VPP (§\,\ref{sec:vpp}). Lui
        donner les départs multiples de \py{robust_solve} est proposé, mais pas fait.
  \item \textbf{Le $\sigma$ du criblage est optimiste} (erreurs supposées indépendantes entre conditions,
        §\,\ref{sec:optimiseur}).
  \item \textbf{Le coût de la sélection.} Le choix d'un lot, en $O(n^2)$, est devenu le goulot : de 25 à
        100 s par lot de 16 au 3\textsuperscript{e} run, pour un plan de 300 à 1300 points. Au
        4\textsuperscript{e}, avec la mise à jour comprise (VPP des 4 sections aux 18 nœuds), un lot passait
        de 1,4 à 3 min au début à environ 6,5 min à la fin, pour un plan de 3\,400 points.
  \item \textbf{Les sensibilités ponctuelles aberrantes} : une fois au 3\textsuperscript{e} run, $d_{F_x}$ a
        valu 201 au lieu de 111 pendant trois mises à jour. On pourrait prendre la médiane de deux solves, ou
        rejeter une valeur trop loin de la précédente.
  \item \textbf{La FSI de premier ordre} et \textbf{la validation entre modèles potentiels} : quelques calculs
        CFD aux états VPP des profils comparés trancheraient. L'étude de maillage du week-end du 23/10
        (\py{runs/2026-10-23_weekend/cfd/report.md}) a montré une traînée convergée mais une portance qui ne
        l'est pas, des couches de prismes mal couvertes et un domaine trop petit. Elle est à refaire avant
        tout calcul de validation.
\end{enumerate}

\end{document}
